%-------------------------
% Resume in Latex
% Author : Harshibar
% Based off of: https://github.com/jakeryang/resume
% License : MIT
%------------------------

\documentclass[letterpaper,11pt]{article}

\usepackage{latexsym}
\usepackage[empty]{fullpage}
\usepackage{titlesec}
\usepackage{marvosym}
\usepackage[usenames,dvipsnames]{color}
\usepackage{verbatim}
\usepackage{enumitem}
\usepackage[hidelinks]{hyperref}
\usepackage{fancyhdr}
\usepackage[english]{babel}
\usepackage{tabularx}
\usepackage{graphicx,scalerel}
\usepackage{fontspec}
% \usepackage[colorlinks = true,
%             linkcolor = blue,
%             urlcolor  = magenta,
%             citecolor = blue,
%             anchorcolor = blue]{hyperref}
% only for pdflatex
% \input{glyphtounicode}
\hypersetup{
    colorlinks=true,
    urlcolor=magenta,
    linkcolor=magenta,
    citecolor=magenta
}
% fontawesome
\usepackage{fontawesome5}

% fixed width
\usepackage[scale=0.90,lf]{FiraMono}

% light-grey
\definecolor{light-grey}{gray}{0.83}
\definecolor{dark-grey}{gray}{0.3}
\definecolor{text-grey}{gray}{.08}

\DeclareRobustCommand{\ebseries}{\fontseries{eb}\selectfont}
\DeclareTextFontCommand{\texteb}{\ebseries}

% custom underilne
\usepackage{contour}
\usepackage[normalem]{ulem}
\renewcommand{\ULdepth}{1.8pt}
\contourlength{0.8pt}
\newcommand{\myuline}[1]{%
  \uline{\phantom{#1}}%
  \llap{\contour{white}{#1}}%
}


% custom font: helvetica-style
\usepackage{tgheros}
\renewcommand*\familydefault{\sfdefault} 
%% Only if the base font of the document is to be sans serif
\usepackage[T1]{fontenc}


\pagestyle{fancy}
\fancyhf{} % clear all header and footer fields
\fancyfoot{}
\renewcommand{\headrulewidth}{0pt}
\renewcommand{\footrulewidth}{0pt}

% Adjust margins
\addtolength{\oddsidemargin}{-0.5in}
\addtolength{\evensidemargin}{0in}
\addtolength{\textwidth}{1in}
\addtolength{\topmargin}{-.5in}
\addtolength{\textheight}{1.0in}

\urlstyle{same}

\raggedbottom
\raggedright
\setlength{\tabcolsep}{0in}

% Sections formatting - serif
% \titleformat{\section}{
%   \vspace{2pt} \scshape \raggedright\large % header section
% }{}{0em}{}[\color{black} \titlerule \vspace{-5pt}]

% TODO EBSERIES
% sans serif sections
\titleformat {\section}{
    \centering \Large \bfseries \vspace{2pt}% \raggedright  % header section
}{}{0em}{}[\color{dark-grey} {\titlerule[2pt]} \vspace{-4pt}]

% only for pdflatex
% Ensure that generate pdf is machine readable/ATS parsable
% \pdfgentounicode=1

%-------------------------
% Custom commands

\newcommand{\resumeItem}[1]{
  \item\small{
    {#1 \vspace{-1pt}}
  }
}

\newcommand{\resumeSubheading}[4]{
  \vspace{-1pt}\item
    \begin{tabular*}{\textwidth}[t]{l@{\extracolsep{\fill}}r}
      \textbf{\textit{#1}} & {\color{dark-grey}\small #2}\vspace{1pt}\\ % top row of resume entry
      \textit{#3} & {\color{dark-grey} \small #4}\\ % second row of resume entry
    \end{tabular*}\vspace{-4pt}
}

\newcommand{\resumeSubSubheading}[2]{
    \item
    \begin{tabular*}{\textwidth}{l@{\extracolsep{\fill}}r}
      \textit{\small#1} & \textit{\small #2} \\
    \end{tabular*}\vspace{-7pt}
}

\newcommand{\resumeProjectheading}[2]{
    \item
    \begin{tabular*}{\textwidth}{l@{\extracolsep{\fill}}r}
      \textbf{\textit{#1}} & {\color{dark-grey} \small #2} \\
    \end{tabular*}\vspace{-4pt}
}

%  ------------MY COMMANDS------------
\newcommand{\sbullet}[1][.5]{\mathbin{\vcenter{\hbox{\scalebox{#1}{$\bullet$}}}}}

\newcommand{\resumeAward}[3]{
    \item[]\hspace*{-\leftmargin}%
    \begin{tabular*}{\dimexpr\linewidth+\leftmargin\relax}{@{}l@{\extracolsep{\fill}}r@{}}\textbf{$\sbullet$\hspace{0.5em}\small #1}, \textit{\small #2} & \small #3 \\
    \end{tabular*}\vspace{-7pt}
}


\newcommand{\resumeLongAward}[3]{
    \item[]
    \begin{tabular*}{0.94\textwidth}{l@{\extracolsep{\fill}}r}
      \textbf{$\sbullet$ \small#1},\\ \textit{\ \ \small #2} & \small #3 \\
    \end{tabular*}\vspace{-7pt}
}




\newcommand{\resumeSubItem}[1]{\resumeItem{#1}\vspace{-4pt}}

% tighter vertical spacing between bullets
\setlist[itemize]{itemsep=0pt, parsep=1pt, topsep=2pt}
\setlist[enumerate]{itemsep=1pt, parsep=0pt, topsep=2pt}

\renewcommand\labelitemii{$\vcenter{\hbox{\tiny$\bullet$}}$}

% CHANGED default leftmargin  0.15 in
\newcommand{\resumeSubHeadingListStart}{\begin{itemize}[leftmargin=0in, label={}]}
\newcommand{\resumeSubHeadingListEnd}{\end{itemize}}
\newcommand{\resumeItemListStart}{\begin{itemize}}
\newcommand{\resumeItemListEnd}{\end{itemize}\vspace{0pt}}
\newcommand{\resumeAwardListStart}{\begin{itemize}[itemsep=4pt, parsep=2pt]}
\newcommand{\resumeAwardListEnd}{\end{itemize}\vspace{0pt}}

\color{text-grey}

%-------------------------------------------
%%%%%%  RESUME STARTS HERE  %%%%%%%%%%%%%%%%%%%%%%%%%%%%


\begin{document}

%----------HEADING----------
\begin{center}
    \textbf{\Huge Chang Liu} \\ \vspace{5pt}
    \small \faEnvelope \hspace{2pt} \texttt{\href{mailto:changl8@andrew.cmu.edu}{changl8@andrew.cmu.edu}} \hspace{1pt} 
    $|$ \hspace{1pt} \faHome \hspace{2pt} \texttt{\href{https://changliu-drpatient.github.io/}{\color{magenta}{changliu-drpatient.github.io}}} \hspace{1pt} 
    \\ \vspace{-3pt}
\end{center}

% \begin{center}
%     Start Date: \textbf{Monday, May 12}; End Date: \textbf{Monday, August 11}
% \end{center}

%-----------EDUCATION-----------
\section{Education}
  \resumeSubHeadingListStart
    \resumeSubheading
      {Kenneth C. Griffin School of Computer Science}{}
      {Carnegie Mellon University, Ph.D. Student in Machine Learning}{08/2023 - 05/2028(est.)}
      \resumeItemListStart
        \resumeItem {Advisor: Prof. Artur Dubrawski.}
        % \resumeItem {QPA: \textbf{4.21/4.33}.}
        \resumeItem {Research Interests: Agentic AI and LLM agents, LLM post-training, systems for AI, machine learning in healthcare.}
        \resumeItem {Courses: Machine Learning for Large Datasets (distributed \& parallel ML), Advanced Deep Learning, Advanced Introduction to Machine Learning, Advanced Machine Learning: Theory and Methods, Intermediate Statistics.}
      \resumeItemListEnd
    \resumeSubheading
      {Yao Class, Institute for Interdisciplinary Information Sciences (IIIS)}{Established by Prof. Andrew C. Yao}
      {Tsinghua University, B.Eng. in Computer Science}{08/2019 - 06/2023}
      \resumeItemListStart
        \resumeItem {GPA: \textbf{3.91/4.00}.}
        \resumeItem {TOEFL: \textbf{120/120}. Reading: 30, Listening: 30, Speaking: 30, Writing: 30.}
        \resumeItem {GRE: \textbf{340/340}. Quant: 170, Verbal: 170, Writing: 5.}
        % \resumeItem {Mathematics Courses: Calculus, Linear Algebra, Abstract Algebra, Mathematics for Computer Science, Mathematics for Artificial Intelligence, 
        %             Probability and Statistics.}
        % \resumeItem {Computer Science Courses: Machine Learning, Reinforcement Learning, Computational Biology, Computer Vision, Deep Learning, Natural Language Processing, Introduction to Databases, 
        %           Data Mining, Quantum Computer Science, Introduction to Robotics, Algorithm Design, Theory of Computation.}
      \resumeItemListEnd
  \resumeSubHeadingListEnd

  \section{Skills}
  \begin{itemize}[leftmargin=0in, label={}, itemsep=0pt, parsep=0pt, topsep=0pt]
    \item \small{\textbf{Languages:} Python, C++, Bash.}
    \item \small{\textbf{AI Agents \& Agentic Systems:} LLM agents, agent flow / agentic workflow design, tool use, multi-agent scaffolds (CodeAct, mini-swe-agent, AWS Strands), prompt engineering.}
    \item \small{\textbf{Systems for AI \& Distributed Computing:} distributed and parallel training (multi-GPU, Ray, FSDP), distributed RL post-training (verl, vLLM), federated learning, large-scale data processing (Spark/MapReduce), efficient LLM inference.}
    \item \small{\textbf{Machine Learning:} LLM post-training (RL, GRPO), retrieval-augmented generation, LLM evaluation and benchmarking, LLM pruning, multimodal ML (EHR, mass spectra), graph neural networks.}
    \item \small{\textbf{Tools:} PyTorch, Hugging Face, Git, Linux.}
  \end{itemize}

  \section{Industry Experience}
  \resumeSubHeadingListStart
    \resumeSubheading
      {TikTok: Video Recommendation}{}
      {Research Intern}{05/2026 - 08/2026}
      \resumeItemListStart
      \resumeItem{Model user preference on PGC (professionally generated content) to boost recommendation performance.}
      \resumeItemListEnd
  \resumeSubHeadingListEnd
  
 
% %-----------EXPERIENCE-----------
% \section{Ongoing Research}
%     \resumeSubHeadingListStart
%           \resumeSubheading
%     {LLM agents for time series machine learning engineering}{07/2025 - Now}
%     {Advised by Prof. Artur Dubrawski.}{Carnegie Mellon University}%\\ \vspace{2pt}
%     \resumeItemListStart
%       \resumeItem {Develop an agentic AI framework for machine learning engineering tasks on time series data.}
%     \resumeItemListEnd
%     \resumeSubheading
%     {Heterogeneous federated learning of foundation models}{03/2025 – Now}
%     {Advised by Prof. Artur Dubrawski.}{Carnegie Mellon University}%\\ \vspace{2pt}
%     \resumeItemListStart
%       \resumeItem {Develop a framework for fine-tuning foundation models with task heterogeneity in a federated learning setting using knowledge distillation.}
%     \resumeItemListEnd
%      \resumeSubheading
%     {Effieicnt Compression of LLMs}{11/2024 – Now}
%     {Advised by Prof. Artur Dubrawski.}{Carnegie Mellon University}%\\ \vspace{2pt}
%     \resumeItemListStart
%       \resumeItem {Explore and innovate depth-based pruning methodologies of LLMs.}
%     \resumeItemListEnd
    
%     \resumeSubHeadingListEnd
\section{Research Experience}
\resumeSubHeadingListStart
    \resumeSubheading
    {Experiential Abstractions in LLMs}{05/2026 - 09/2026}
    {LLM post-training. Distributed RL training (verl). Retrieval-augmented LLMs.}{Carnegie Mellon University}%\\ \vspace{2pt}
    \resumeItemListStart
      \resumeItem {Develop a novel framework for improving LLM problem-solving abilities by first extracting abstractions from training, and then absorbing them via in-context learning or RL.}
      \resumeItem {Run multi-GPU distributed GRPO training with verl (Ray-orchestrated FSDP training workers and vLLM rollout engines).}
      \resumeItem {Improve mathematical reasoning abilities compared to naive GRPO, even without a strong teacher for abstraction extraction. pass@1/pass@8: \textbf{46.89/70.54 $\to$ 49.07/71.14}. }
      \resumeItem{Accepted at EMNLP Findings. Paper link \href{https://arxiv.org/abs/2607.20372}{here}.}
    \resumeItemListEnd
    
    \resumeSubheading
    {Lightweight LLM agents for ML engineering (Agentic AI)}{05/2025 - 08/2025}
    {LLM agents. Agent flow design. Tool use. AWS Strands. Scaffold engineering.}{Carnegie Mellon University}%\\ \vspace{2pt}
    \resumeItemListStart
        \resumeItem {Develop a lightweight agentic system that autonomously completes end-to-end machine learning engineering (MLE) tasks.}
      \resumeItem {Design a multi-step agent flow following the CodeAct paradigm (draft, execute, revise), delegating code edits to a \href{https://github.com/swe-agent/mini-swe-agent}{mini-swe-agent} sub-agent invoked as a tool.}
      \resumeItem {Use AWS's \href{https://strandsagents.com/}{strands} library for building the agents. Project link \href{https://github.com/ChangLiu-DrPatient/miniMLEagent}{here}. }
    \resumeItemListEnd
    
    \resumeSubheading
    {Depth pruning and emergent abilities of pruned LLMs}{12/2024 - 09/2025}
    {System for AI. Efficient LLM inference. LLM pruning. Hugging Face. Benchmarking.}{Carnegie Mellon University}%\\ \vspace{2pt}
    \resumeItemListStart
      \resumeItem {Propose a novel depth-based LLM pruning method that retains key reasoning abilities under \textbf{10}-fold compression.}
      \resumeItem{Systematically benchmarked the emergent abilities of aggressively pruned LLMs. Paper link \href{https://openreview.net/forum?id=aP3sv2xrGH}{here}.}
    \resumeItemListEnd
    % \resumeSubHeadingListEnd


% \section{Healthcare Research Experience}
%     \resumeSubHeadingListStart
    \resumeSubheading
    {Multimodal outbreak detection in hospitals}{10/2023 - Now}
    {Multimodal ML. ML for healthcare. }{Carnegie Mellon University}%\\ \vspace{2pt}
    \resumeItemListStart
      % \resumeItem {Developed a semi-supervised pretraining and finetuning framework with a UNet-based autoencoder for MALDI-TOF data.}
      \resumeItem {Develop a multimodal machine learning framework for hospital outbreak detection that integrates signals from MALDI-TOF mass spectra, antimicrobial resistant patterns, and electronic health records (EHR).}
      \resumeItem {Enable cost-effective outbreak detection that can reduce dependency on whole genome sequencing gold standard by \textbf{80\%} while retaining \textbf{80\%} discriminatory power.}
    \resumeItemListEnd


    \resumeSubheading
      {Multimodal disease target identification through a probabilistic knowledge graph}{09/2021 - 05/2023}
      {Multimodal ML. Graph Neural Networks. Computational biology.}{Tsinghua University}%\\ \vspace{2pt}
      \resumeItemListStart
        \resumeItem {Develop a novel method of augmenting biological networks with literature evidence to construct a probabilistic knowledge graph.} 
        \resumeItem {Develop a graph neural network to predict target candidates from the knowledge graph, achieving superior performance to state-of-the-art 
        models in terms of accuracy (esp. on sparse data) and literature support for top novel predictions.}
        \resumeItem {Conduct bioinformatics analyses and cooperate with experimental validation of the identified colorectal cancer and melanoma targets.}
      \resumeItemListEnd

      % \resumeSubheading
      % {Reconstructing the Allele-specific Genome Structure from Hi-C Contacts}{03/2022 - 03/2023}
      % {Advised by Prof. Jian Ma.}{Carnegie Mellon University}%\\ \vspace{2pt}
      % \resumeItemListStart
      %   \resumeItem {Developed an improved particle dynamics framework (based on \textit{hickit})
      %   that iterates between inferring chromosome contact phases and 3D genomic coordinates to fully exploit their common information.}
      %   \resumeItem {Developed a new graph neural network to implicitly impute the phases of the 
      %   Hi-C contacts and reconstruct the allele-specific 3D genome structure.}
      % \resumeItemListEnd

    % \resumeSubheading
    %   {Discovering Competitive Binding of Transcription Factors}{05/2021 - 02/2023}
    %   {Advised by Prof. Jianyang Zeng.}{Tsinghua University}
    %   \resumeItemListStart
    %     \resumeItem {Developed a framework to infer in-vivo competitive TF binding 
    %     (the binding of one TF removes that of the other), consisting of a deep neural network, several motif analyses, and statistical tests.}
    %     \resumeItem {Cooperated with experimental validation of the predicted competing TF pairs (in progress).}
    %   \resumeItemListEnd
    
    % \resumeSubheading
    %   {Predicting Antigen Binding Sites through Graph Neural Networks}{06/2021 - 08/2021}
    %   {Advised by Prof. Boxue Tian.}{Tsinghua University}
    %   \resumeItemListStart
    %     \resumeItem {Developed a graph neural network to predict antigen binding residues using 
    %     antigen-antibody compound data in the SAbDAb database based on \href{http://www.csbio.sjtu.edu.cn/bioinf/GraphBind/sourcecode.html}{\textit{GraphBind}}, 
    %     a DNA/RNA-Protein binding site prediction model.} 
    %     \resumeItem {Utilized the model to validate lab-generated compounds.}
    %   \resumeItemListEnd

    % \resumeSubheading
    %   {Intelligent Diabetes Management}{12/2020 - 02/2021}
    %   {Advised by Prof. Yang Yuan.}{Tsinghua University}
    %   \resumeItemListStart
    %     \resumeItem {Cooperated with Shanghai Zhongshan Hospital to investigate the needs of the endocrinology department and its patients.}
    %     \resumeItem {Developed a deep learning framework for predicting future patient blood sugar levels from patient records for pre-emptive alerts.} 
    %     \resumeItem {Developed a deep learning framework for predicting the proper dosage of insulin to be administered to alleviate the demand for 
    %     expert consultation.}
    %   \resumeItemListEnd

    % \resumeProjectheading
    %   {Course Projects}{Tsinghua University}
    %   \resumeItemListStart
    %     \resumeItem {Separating Bone Inscription Images through Computer Vision Techniques}
    %     \resumeItem {Avoiding Catastrophic Failure in Reinforcement Learning}
    %     \resumeItem {Encoding Semantics into Word Embeddings through Dictionaries}
    %     \resumeItem {Detecting and Ranking Communities in the DBLP Database}
    %   \resumeItemListEnd
  \resumeSubHeadingListEnd




%
%-----------PROGRAMMING SKILLS-----------

\section{Selected Publications}
\begin{enumerate}
    \item Alexander J. Sundermann$^*$; Jieshi Chen; \textbf{Chang Liu};...; Artur Dubrawski; Lora Lee Pless; Lee H. Harrison. ``MALDI-TOF for healthcare transmission prediction: comparison with whole genome sequencing'' \textit{Infection Control \& Hospital Epidemiology}, October 2026.
    
    \item \textbf{Chang Liu}; Xinyu Li; Artur Dubrawski$^*$.``Notes to Self: Can LLMs Benefit from Experiential Abstractions?''  \textit{EMNLP Findings}, August 2026. paper link \href{https://arxiv.org/abs/2607.20372}{here}.
    
    \item \textbf{Chang Liu}; Jieshi Chen; Alexander J. Sundermann; Kathleen Shutt; Marissa P. Griffith; Lora Lee Pless; Lee H. Harrison; Artur Dubrawski$^*$.``Towards Practical Multimodal Hospital Outbreak Detection,'' \textit{AMIA Annual Symposium}, May 2026.
        
    \item \textbf{Chang Liu}; Jieshi Chen; Alexander J. Sundermann; Kathleen Shutt; Marissa P. Griffith; Lora Lee Pless; Lee H. Harrison; Artur Dubrawski$^*$.``Exploring the Utility of MALDI-TOF Mass Spectrometry and Antimicrobial Resistance in Hospital Outbreak Detection,'' \textit{AMIA Amplify Informatics Conference}, February 2026.

    
    \item \textbf{Chang Liu$^\dag$}; Arjun Choudhry$^\dag$; Yifu Cai; Nina \.{Z}ukowska; Mononito Goswami; Artur Dubrawski. ``Depth as a Scaling Vector: Simple Pruning and Evaluation of Emergent Abilities in Pruned LLMs,'' \textit{NeurIPS 2025 Workshop on Evaluating the Evolving LLM Lifecycle: Benchmarks, Emergent Abilities, and Scaling}, September 2025.

    % \item Arjun Choudhry$^\dag$; \textbf{Chang Liu$^\dag$}; Nina \.{Z}ukowska; Yifu Cai; Mononito Goswami; Artur Dubrawski. ``LayerMerge: Modality-Agnostic Depth Pruning for Efficient Foundation Model Deployment,'' \textit{NeurIPS 2025 Workshop on Efficient Reasoning}, September 2025.

    % \item Miko{\l}aj Pi\'{o}rczy\'{n}ski; Wojciech {\L}apacz; Xinyu Li; \textbf{Chang Liu}; Abby Turner; Artur Dubrawski. ``From Aggregation to Guidance: Strategies for Personalized Federated Fine-Tuning of Foundation Models，'' \textit{NeurIPS 2025 Workshop on Unifying Representations in Neural Models}, September 2025.
    
    \item \textbf{Chang Liu}; Jieshi Chen; Lee H. Harrison; Artur Dubrawski$^*$. ``Bridging the Utility Gap Between MALDI-TOF and WGS for Affordable Outbreak Cluster Detection,'' \textit{The AHLI Conference on Health, Inference, and Learning (CHIL)}, April 2025.
    
    % \item \textbf{Chang Liu}; Jieshi Chen; Lee H. Harrison; Artur Dubrawski$^*$. ``Multimodal Structure Preservation Learning,''\textit{arXiv preprint}, October 2024.
  
    \item \textbf{Chang Liu$^\dag$}; Kaimin Xiao$^\dag$; Cuinan Yu$^\dag$; Yipin Lei$^\dag$;...;
    Dan Zhao$^*$; Fengfeng Zhou$^*$; Haidong Tang$^*$; Jianyang Zeng$^*$. 
  ``A Probabilistic Knowledge Graph for Target Identification,'' \textit{PLOS Computational Biology}, April 2024.
  
  % \item \textbf{Chang Liu$^\dag$}; Cuinan Yu$^\dag$; Yipin Lei$^\dag$;...;
  % Dan Zhao$^*$; Fengfeng Zhou$^*$; Jianyang Zeng$^*$.
  % ``Improving Target-disease Association Prediction through a Graph Neural Network with Credibility Information,''
  % proceedings of the \textit{Pacific Symposium on Biocomputing}, January 2023. \\
\end{enumerate}

\section{Honors \& Awards}
 \resumeAwardListStart
  \resumeAward{Comprehensive Merit Award (7/32)}{Tsinghua University}{2022}
% \resumeAward{3rd place in the 100m Men's Freestyle}{Tsinghua University}{2021}
  \resumeAward{Comprehensive Merit Award (6/32)}{Tsinghua University}{2021}
  % \resumeLongAward{Gold Prize in the Magnus Royal Youth International Piano Competition Semi-Final in China}{Magnus Conservatory of Music, Lithuania}{2021}
  % \resumeLongAward{2nd place in the National College Water Polo Championship (university water polo team)}{Federation of University Sports of China}{2021}
  % \resumeAward{4th place in the 50m Men's Freestyle}{Tsinghua University}{2021}
  % \resumeAward{1st place in the 200m Mens' Medley}{Tsinghua University}{2021}
  % \resumeAward{37th place in the Campus Mini Marathon}{Tsinghua University}{2020}
  % \resumeAward{Excellence Award for Volunteering Services}{Tsinghua University}{2020}
  % \resumeAward{Freshmen Scholarship}{Tsinghua University}{2019}
  % \resumeAward{University Full Scholarship for Future Scholars}{Tsinghua University}{2019}
 \resumeAwardListEnd

% \section{Academic Service}
% \resumeItemListStart
% \resumeItem{\textbf{Peer Review:} Reviewed 19 academic papers. NeurIPS Neural Artifacts Workshop 2026 (3), NeurIPS 2026 (1), NeurIPS UniReps Workshop 2025 (2), AMIA Amplify Informatics Conference 2026 (5), AMIA Annual Symposium (8).}
% \resumeItemListEnd
%-------------------------------------------
\end{document}
